% Appendix D

\chapter{Configuración de los modelos} % Main appendix title

\label{AppendixD}

En este apéndice se presentan el espacio de búsqueda de los hiperparámetros de LGBM\_OPT y los umbrales de los grados-día, descritos en la sección~\ref{sec:desarrollo_modelos}.

%----------------------------------------------------------------------------------------
\section{Espacio de búsqueda de los hiperparámetros}
\label{sec:apendice_optuna}
%----------------------------------------------------------------------------------------

La tabla~\ref{tab:espacio_optuna} presenta los rangos explorados con Optuna en la búsqueda de los hiperparámetros de LGBM\_OPT, descrita en la subsección~\ref{subsec:modelos_profundos}.

\begin{table}[H]
	\centering
	\caption{Espacio de búsqueda de los hiperparámetros de LGBM\_OPT.}
	\label{tab:espacio_optuna}
	\footnotesize
	\setlength{\tabcolsep}{4pt}
	\renewcommand{\arraystretch}{1.15}
	\begin{tabularx}{\textwidth}{
			@{}
			>{\raggedright\arraybackslash\hsize=1.2\hsize}X
			>{\raggedright\arraybackslash\hsize=0.8\hsize}X
			@{}
		}
		\toprule
		\textbf{Hiperparámetro} & \textbf{Rango} \\
		\midrule
		Cantidad de árboles & 100 a 1000, en pasos de 50. \\
		Tasa de aprendizaje & 0,01 a 0,2, escala logarítmica. \\
		Cantidad de hojas & 8 a 128, escala logarítmica. \\
		Mínimo de observaciones por hoja & 5 a 100, escala logarítmica. \\
		Fracción de observaciones por árbol & 0,5 a 1. \\
		Fracción de variables por árbol & 0,5 a 1. \\
		Regularización L2 & $10^{-3}$ a 10, escala logarítmica. \\
		\bottomrule
	\end{tabularx}
\end{table}

%----------------------------------------------------------------------------------------
\section{Umbrales de los grados-día}
\label{sec:apendice_umbrales}
%----------------------------------------------------------------------------------------

La tabla~\ref{tab:umbrales_grados_dia} presenta los umbrales de grados-día, descritos en la subsección~\ref{subsec:respuesta_temperatura}.

\begin{table}[H]
	\centering
	\caption{Umbrales de grados-día por región eléctrica, en \textdegree C.}
	\label{tab:umbrales_grados_dia}
	\footnotesize
	\setlength{\tabcolsep}{4pt}
	\renewcommand{\arraystretch}{1.15}
	\begin{tabular}{@{}lcccc@{}}
		\toprule
		& \multicolumn{2}{c}{\textbf{Temperatura aparente}} & \multicolumn{2}{c}{\textbf{Temperatura del aire}} \\
		\cmidrule(lr){2-3} \cmidrule(l){4-5}
		\textbf{Región} & \textbf{Refrigeración} & \textbf{Calefacción} & \textbf{Refrigeración} & \textbf{Calefacción} \\
		\midrule
		CEN & 20,0 & 12,5 & 20,0 & 15,5 \\
		COM & 17,0 & 12,5 & 19,5 & 15,5 \\
		CUY & 17,5 & 13,0 & 20,5 & 15,5 \\
		GBA & 23,0 & 17,0 & 22,0 & 18,0 \\
		LIT & 21,0 & 14,0 & 21,0 & 16,0 \\
		NEA & 22,0 & 12,5 & 21,0 & 15,0 \\
		NOA & 21,0 & 14,5 & 21,0 & 15,5 \\
		PAT & 14,5 & 14,5 & 19,0 & 17,0 \\
		PBN & 21,5 & 15,5 & 21,0 & 18,5 \\
		PBS & 17,0 & 12,0 & 19,5 & 16,0 \\
		\bottomrule
	\end{tabular}
\end{table}
